\hypertarget{appendix-dbe5905163c2}{}
\section*{Bit mask}
\section{Bit mask}\label{bit-mask}

\textbf{Curriculum depth:} Introduced · \textbf{First seen:} C001 · \textbf{Latest development:} C001

\subsection{Definition and mechanism}\label{definition-and-mechanism}

A pattern naming selected bit positions; its effect depends on the operation used with it.

A mask is a selection pattern, not an operation. \texttt{0x30\ =\ 00110000₂} names positions 4 and 5. AND reads selected positions, OR sets them, XOR toggles them, and AND with the width-limited complement clears them. Full-width masks need a special implementation case.

\subsection{Prerequisites}\label{prerequisites}

\begin{itemize}
\tightlist
\item
  \hyperref[appendix-280445746493]{Bit significance} --- LSB means least significant bit, normally position 0 with weight 1.
\item
  \hyperref[appendix-5bee2df85484]{Byte and octet} --- An octet is exactly eight bits.
\end{itemize}

\subsection{Used by}\label{used-by}

\begin{itemize}
\item
  \hyperref[appendix-ac697f5c66e8]{Bit field} --- uses bit mask as part of its representation or reasoning.
\item
  \hyperref[appendix-0c67cd5bf30a]{Pages and page offsets} --- depends on this concept; see its prerequisite explanation.
\end{itemize}

\subsection{Recall and next depth}\label{recall-and-next-depth}

Explain bit mask without the label, work the example above, then state the assumption whose removal would change the result.

\textbf{Pending:} Apply and reconstruct this concept independently.

\subsection{Source and navigation}\label{source-and-navigation}

\href{Continuous\%20Notes/Computing\%20Foundations\%20-\%20From\%20Bits\%20to\%20C++\%20and\%20Systems\%20APIs\#From\%20a\%20signal\%20to\%20an\%20interpretation}{C001 teaching chapter} · \hyperref[appendix-fd090d982a15]{Part definitions and mechanisms} · \href{}{Dependency Map (02-Foundations-II, appendix-4ad47bdd269a)}
